\section{为什么需要这座桥}

本章把产业链逻辑转成盈利逻辑。AI平台升级本身不是投资结论，只有当平台需求能落到公司收入、毛利率、订单持续性和EPS弹性上，才构成可投资判断。

\begin{houseviewbox}[Research upgrade required]
顶级半导体硬件研报不会停留在``AI平台升级带来PCB需求''，而会把平台链条转成公司收入、毛利和EPS弹性。本章明确给出当前可用的客户链桥，并把缺失变量标为 corpus gap。
\end{houseviewbox}

客户链是本报告最重要的分析轴。AI服务器带来的不是均匀增长，而是订单向少数能通过客户认证、能承接高层数和低损耗材料要求的供应商集中。因此，判断一家公司是否真正受益，不能只看“是否在AI PCB概念里”，而要看三件事：是否有可验证的应用收入，是否有客户或终端认证线索，是否能把收入增长转成毛利率和现金流。

这也是为什么同样处在PCB产业链中，沪电股份、胜宏科技、深南电路、生益科技、华正新材的投资含义完全不同。沪电股份更接近高端数通PCB的确定性资产；胜宏科技更像高弹性客户链资产；深南电路有PCB和载板双轮驱动，但载板良率与折旧是利润转换变量；生益科技的核心不是PCB产能，而是M8/M9高频高速材料的价格和认证；华正新材则更多是国产材料期权，必须用订单和利润修复来验证。

\section{平台到盈利桥}

先看平台到盈利的路径。NVIDIA/Rubin、ASIC/TPU、国产算力、光模块和先进封装都会带来PCB/CCL需求，但每一条链条对应的缺口不同：有的缺客户份额，有的缺单机价值量，有的缺良率和折旧桥。表格的作用是说明“哪些链条能进入模型，哪些仍然只能作为情景变量”。

需求侧公开证据继续强化。三大云厂2026Q1收入与CAPEX均创新高，且均明确指向AI基础设施。\textbf{Alphabet}：2026Q1 Google Cloud收入同比增长63\%至200亿美元，由企业AI解决方案和AI基础设施拉动，Q1购置物业和设备356.74亿美元。\textbf{Amazon}：2026Q1 AWS收入同比增长28\%至376亿美元，自由现金流下降主要来自物业和设备购置同比增加593亿美元、且主要反映AI投资，并披露Trainium、NVIDIA GPU和AI芯片部署。\textbf{Meta}：2026Q1资本开支198亿美元，由服务器、数据中心和网络基础设施驱动，并将2026全年资本开支指引上调至1250--1450亿美元。上述证据可以证明AI基础设施CAPEX仍是需求驱动，但不能映射为任何PCB/CCL公司的具名客户收入、订单、ASP或EPS贡献。

\begin{exhibitbox}[图表5-1: Platform-chain bridge and corpus gaps]
\footnotesize
\begin{tabularx}{\textwidth}{L{2.6cm}X X X}
\toprule
\textbf{Platform chain} & \textbf{PCB / CCL implication} & \textbf{A-share names} & \textbf{Missing earnings variable} \\
\midrule
NVIDIA GB300 / Rubin & 高层AI PCB、背板/中板、M9材料、CoWoP相关结构。 & 沪电股份、胜宏科技、深南电路、生益科技 & 客户份额、单机价值量、订单拆分、ASP。 \\
North American ASIC / TPU & HDI、高速服务器PCB、交换机板、M8/M9 CCL。 & 胜宏科技、沪电股份、生益科技 & 项目收入、客户集中度、预测明细。 \\
Domestic compute & 国产高端CCL、CBF/BT、PCB与载板认证。 & 华正新材、深南电路、生益科技、兴森科技 & 客户验证、放量节奏、价格与良率。 \\
Networking / optical & 光模块PCB、交换机PCB、高速材料。 & 沪电股份、深南电路、景旺电子等 & 平台 mix、光模块/交换机收入贡献。 \\
Advanced packaging / substrate & FC-BGA、ABF/BT、玻璃基板、CoWoP-like convergence。 & 深南电路、兴森科技、华正新材 & 良率、产能、折旧、客户认证。 \\
\bottomrule
\end{tabularx}
\sourcenote{house-view customer-chain bridge; hyperscaler CAPEX evidence (official Q1 releases); hyperscaler capex supports demand bridge only, not supplier revenue allocation}
\end{exhibitbox}

表中最关键的结论是：目前公开材料能够证明AI平台是需求来源，但不能完整证明每个平台对每家公司EPS的贡献。因此，对核心公司可以给更高跟踪优先级，但对具名NVIDIA、Google ASIC、Apple或国产算力平台的收入拆分，仍必须保持证据边界，不能把市场传闻直接写进估值模型。

\begin{exhibitbox}[图表5-2: Named platform revenue split confirmation pack]
\footnotesize
\begin{tabularx}{\textwidth}{L{2.0cm}>{\raggedright\arraybackslash}X >{\raggedright\arraybackslash}X >{\raggedright\arraybackslash}X}
\toprule
\textbf{公司} & \textbf{已有公开收入基线} & \textbf{需要直接确认的拆分} & \textbf{当前披露边界} \\
\midrule
沪电股份 & 2025 PCB收入约181.43亿元；数通PCB约146.56亿元；高速交换机/路由81.69亿元；AI服务器/HPC 30.06亿元。 & NVIDIA/Rubin/GB系列、Google TPU/ASIC、AMD/海外云、国产算力、高速交换机/光模块平台收入。 & 互动易回复称具体厂商和客户合作受商业政策限制，不便讨论。 \\
胜宏科技 & 2025收入192.92亿元；官方材料确认高密度PCB、100层以上高多层板、HDI和AI Data Center UBB/交换机能力；互动易确认ASIC业务推进和1.6T光模块mSAP订单饱满。 & NVIDIA系（Blackwell/Rubin/GB200/GB300）；Google TPU/ASIC；Microsoft/Amazon ASIC；UBB/OAM/midplane/正交背板；1.6T光模块mSAP收入。 & 互动易回复称客户名称和业务细节未经许可不能讨论。 \\
深南电路 & 2025 PCB收入143.59亿元、载板收入41.48亿元；IR确认高速交换机、光模块、AI加速卡和服务器需求。 & AI加速卡PCB、高速交换机、光模块PCB、服务器/存储PCB、FC-BGA/ABF/BT载板客户或认证状态。 & P5W/深交所回复称涉及具体企业合作的问题基于商业保密不便回复。 \\
生益科技 & 官方与IR证据支持高频高速材料、封装基板材料，以及国内外终端GPU/AI项目合作。 & M8/M9/M10收入、GPU AI服务器材料、ASIC AI服务器材料、国产算力、交换机/光模块材料收入和具名认证状态。 & 上证e互动IR不披露具体客户或M8/M9/M10收入占比；问询材料对客户名称和信用政策豁免披露。 \\
华正新材 & 2025收入43.69亿元，CCL收入33.75亿元；官方文本指向AI服务器、交换机、光模块、天线及CBF/BT/高等级CCL。 & 高速CCL按AI服务器/交换机/光模块/天线拆分，BT/类BT、CBF/CBF-RCC按客户/应用拆分，国产算力验证和爬坡收入。 & 官方文件和互动平台只有产品/应用进度，不披露具名客户收入、ASP、出货或订单金额。 \\
\bottomrule
\end{tabularx}
\sourcenote{named-platform revenue split confirmation pack; official filing segment summary; customer disclosure boundary evidence; direct confirmation pack, not completed revenue evidence}
\end{exhibitbox}

这个确认包把剩余工作从“继续泛泛找客户链线索”改成“向公司IR、客户/供应商确认或付费数据库索取指定字段”。只有当可归因来源给出客户/平台收入拆分，或公司/客户/供应商明确确认该颗粒度不可披露时，具名平台收入拆分这一项才能关闭。

\section{公司桥}

把平台链条落到公司层面后，差异更明显。沪电股份已有AI服务器/HPC和高速交换机收入线索，说明它不是纯概念；胜宏科技的能力和客户弹性更强，但对客户集中度和A/H估值可比性更敏感；深南电路的优势是财务交付稳健，但IC载板和FC-BGA的盈利质量仍需持续验证；生益科技要看高端CCL价格和材料代际；华正新材要看CBF/BT小批量订单能否形成可重复收入。

\begin{exhibitbox}[图表5-3: Company bridge from chain exposure to verification]
\footnotesize
\begin{tabularx}{\textwidth}{L{2.2cm}X X X}
\toprule
\textbf{公司} & \textbf{Observable exposure} & \textbf{Direct financial evidence} & \textbf{Upgrade evidence needed} \\
\midrule
沪电股份 & AI PCB、backplane、CoWoP相关需求。 & Goldman目标价142元（基于23x 2027E EPS，Sina转载非原始PDF，基准约2026-05），管理层调研记录。 & AI收入占比、平台 mix、ASP/毛利、EPS明细。 \\
胜宏科技 & NVIDIA升级周期、Google ASIC/TPU。 & JPM给出2025--2028净利CAGR与H股目标价约600港元（财经媒体转载非原始PDF，基准约2026-05）。 & 原始报告、客户集中度、产能爬坡、A/H可比性。 \\
深南电路 & AI算力PCB + IC载板。 & 2025、2026Q1和PCB/载板分部数据；2026--2028预测。 & 客户平台拆分、FC-BGA/BT良率、载板折旧。 \\
生益科技 & AI CCL、M9迁移、材料紧缺。 & Goldman目标价217.6元（基于CCL涨价逻辑，Sina转载非原始PDF，基准约2026-05）。 & M8/M9收入占比、涨价幅度、客户认证。 \\
华正新材 & 国产算力CCL、CBF/BT材料。 & 海通国际和浙商2023 CBF/高速CCL深度、2026深度和年报点评补充CBF技术路径、2023E--2028E收入/净利预测，官方年报补充客户/供应商集中度。 & CBF/BT客户验证、产品良率、客户/平台拆分和订单转化。 \\
\bottomrule
\end{tabularx}
\sourcenote{company bridge; Goldman (沪电142元、生益217.6元) / JPM (胜宏H股约600港元) target-price and CAPEX reprints via Sina / financial media, not original broker PDFs; benchmark approx 2026-05; segment data from official filings}
\end{exhibitbox}

\section{完整报告中的量化桥}

\paragraph{本章统一证据边界} 在进入逐公司证据之前，先界定本章统一的披露边界：已公开的分部收入、毛利率、客户/供应商集中度数据可作为模型输入，但\textbf{具名平台（NVIDIA/Rubin、Google TPU/ASIC、Microsoft/Amazon ASIC、光模块等）收入拆分、单机价值量（ASP）和平台份额一律未由任何官方或一手文件披露}，仅能作为待验证情景变量。下文每个公司的``缺口''项只标注该公司特有的未披露内容，不再重复上述通用边界。

进一步看量化桥，报告能直接使用的硬证据主要来自年报、H股申请资料、IR记录和已归档券商模型。这里的重点不是简单罗列每家公司有什么资料，而是识别哪些资料能够进入盈利模型。收入分部、毛利率、产能利用率和订单储备可以进入模型；具体客户名称、单机价值量和平台份额如果没有公开披露，只能作为待验证变量。

\par\vspace{6pt plus 2pt minus 2pt}
\footnotesize
\begin{longtable}{@{}L{1.6cm}p{7.6cm}p{6.8cm}@{}}
\multicolumn{3}{@{}l@{}}{\textbf{\color{navy}图表5-4: Quantified platform-to-company bridge from downloaded PDFs}} \\
\noalign{\vskip 3pt \color{navy}\hrule height 1.2pt \color{black}\vskip 5pt}
\textbf{公司} & \textbf{硬证据} & \textbf{缺口} \\
\midrule
\endfirsthead
\multicolumn{3}{@{}l@{}}{\textit{图表5-4（续）：Quantified platform-to-company bridge}} \\
\toprule
\textbf{公司} & \textbf{硬证据} & \textbf{缺口} \\
\midrule
\endhead
\midrule
\multicolumn{3}{@{}r@{}}{\textit{续下一页}} \\
\endfoot
\bottomrule
\noalign{\vskip 3pt \color{rulegrey}\hrule height 0.35pt \color{black}}
\endlastfoot
沪电股份 & 2025数通PCB收入146.56亿元，高速交换机81.69亿元，AI服务器/HPC 30.06亿元；2026Q1收入62.14亿元、净利12.42亿元。注：官方分部口径下，高速交换机81.69+AI服务器/HPC 30.06+通用服务器25.40+无线/网络其他9.41=146.56亿元，与数通PCB分部精确吻合，故四项为数通PCB子集；但AI服务器/HPC 30.06亿元内部仍无法拆出NVIDIA/Google/国产算力平台占比。2026年6月IR补充泰国数通事业部海外客户认证超70\%、产能利用率超90\%，并加速AI服务器与高速网络产品导入交付。CNInfo/深交所互动易原始接口显示，公司对Rubin、谷歌、NVIDIA/AMD/海外云厂商供货比例、LPU/LPX和国产算力客户等具体厂商问题，均以商业政策限制为由不予讨论；2026年6月18日互动易局部刷新进一步补充高端PCB定价随技术价值、质量溢价和供需动态调整，CoWoP/mSAP/光铜融合/M10仍有研发周期和商业化风险，泰国基地正从爬坡期迈入规模化运营期。 & AI客户份额和单机价值量仍未披露。 \\
\midrule
胜宏科技 & 2026Q1收入55.19亿元、归母净利12.88亿元；2026年6月IR披露具备100层以上高多层PCB、10阶30层HDI、16层任意互联HDI能力，下一代14阶36层HDI推进研发认证；CNInfo/深交所互动易原始接口确认ASIC相关客户业务进展顺利、1.6T光模块mSAP产能用于满足订单且在手订单饱满，同时对NVIDIA Spark、midplane、mSAP/正交背板、CB300、国内芯片、Tesla AI5、GB200/GB300和CoWoP等具体客户/项目问题以商业政策限制为由不披露；局部刷新确认惠州厂房四分阶段投产、厂房十/十一及泰国工厂建设按计划推进。 & AI数据中心UBB/交换机客户收入占比未披露。 \\
\midrule
深南电路 & 2025 PCB收入143.59亿元、毛利率35.53\%；封装基板收入41.48亿元、毛利率22.58\%（以上为公司分部披露）；券商估算（CMBI，非公司披露）称AI data center约占PCB收入25\%；2026年6月IR确认PCB业务受AI算力基础设施硬件需求拉动、产能利用率维持高位，FC-BGA 22层及以下量产、24层以上研发打样推进。P5W/深交所互动易原始接口进一步确认，公司承认高速通信网络、数据中心交换机、AI加速卡、存储等领域需求受AI趋势影响，但对具体企业合作问题以商业保密原则不予回复。 & AI服务器订单金额和FC-BGA客户结构未披露。 \\
\midrule
生益科技 & 2026Q1毛利率28.10\%，同比+3.50pct、环比+2.32pct；2025H1官方分部披露覆铜板外部收入84.67亿元、PCB外部收入37.53亿元；生益电子问询回复披露5--6阶HDI已获境内外知名客户认证并小批量销售；上证e互动探测显示近期M8/M9相关提问存在，但回复流未给出新的发行人答复，官方IR PDF仍是有效证据边界。 & M8/M9收入占比、涨价幅度、具名客户和客户收入仍缺。 \\
\midrule
华正新材 & 海通国际2023年13页原始报告补充高速CCL由M6向M7/M8升级、56Gbps数通交换机小批量、400G光模块/高阶AI服务器新序列开发和43元历史目标价；浙商2023年17页CBF深度报告补充FC-BGA等先进封装材料原理、日企约96\%份额和2023E--2025E分产品收入/毛利率模型；浙商41页深度报告及年报点评补充高端CCL、CBF膜、BT封装材料和2025E--2028E模型；官方可行性报告补充1200万张高等级覆铜板项目；P5W/上交所互动显示CBF部分单品小批量订单销售、其他单品终端验证推进，类BT/BT板材已在内存封装等领域应用。 & 具名客户、CBF/BT收入和客户/平台收入拆分仍缺。 \\
\end{longtable}
\sourcenote{full report model summary; supplemental report archive summary; official IR refresh (2026年6月); customer disclosure boundary evidence; CNInfo/SZSE/SSE interaction sweep; Shengyi SSE interaction probe; company filings / IR records / issuer interaction evidence / full reports; customer-specific revenue remains unavailable}
\par\vspace{6pt plus 2pt minus 2pt}
\normalsize


量化桥给出的结论是偏谨慎的。沪电股份和胜宏科技已经有足够多的高端PCB应用证据，但仍无法拆分具体平台收入；深南电路的财务兑现更稳，但载板的利润转换还没有完全证明；生益科技的材料逻辑强，但M8/M9收入占比和涨价幅度仍是关键空白；华正新材最需要从“材料期权”转向“订单和利润证明”。这决定了后续估值不能一刀切地套用AI硬件高倍数。

\begin{exhibitbox}[图表5-5: H-share prospectus customer evidence]
\footnotesize
\begin{description}
  \item[沪电股份] H股申请资料披露AI服务器/HPC收入：2023年12.40亿元/13.9\%，2024年29.75亿元/22.3\%，2025年30.06亿元/15.9\%，2026Q1为7.34亿元/18.2\%；前五大客户收入占比从2023年的46.0\%升至2026Q1的58.4\%，最大客户占比升至19.5\%。
  \item[胜宏科技] H股全球发售资料披露客户类别包括全球AI技术解决方案供应商、大型云服务商、数据中心设备OEM、服务器厂商等；前五大客户收入占比2023/2024/2025年为27.1\%/25.1\%/51.0\%，最大客户占比为6.3\%/8.3\%/29.7\%，海外收入占比2025年达76.8\%。
  \item[客户质量] 沪电前五大客户主要销售数据通讯PCB，信用期30--120天或75--110天，多个客户合作起始于2001--2018年；胜宏Customer A在2025年收入57.38亿元/29.7\%，胜宏Customer E被描述为总部位于美国、纳斯达克上市、聚焦加速计算和AI基础设施的全球领先科技公司（注：``Customer E''为胜宏H股招股书的匿名客户代码，与鹏鼎H股文件中的``Customer E''分属不同公司、不同实体，两份招股书的代码体系无对应关系，不可合并解读为同一终端客户）。
  \item[边界] 两份H股文件显著增强客户集中度和应用收入证据，但仍未披露NVIDIA/Google ASIC/光模块等具名平台收入。
\end{description}
\sourcenote{HKEX H-share application proof and global offering prospectus}
\end{exhibitbox}

H股文件强化了客户集中度判断：AI订单越强，客户集中度往往越高，这既是增长来源，也是估值折扣来源。对胜宏科技和鹏鼎控股尤其如此，客户集中度可以解释盈利弹性，也意味着一旦大客户节奏变化，利润和估值会同步受压。一个值得注意的反直觉信号是：沪电AI服务器/HPC收入绝对值从2024年29.75亿元增至2025年30.06亿元，但占比却从22.3\%回落至15.9\%，反映数通PCB整体放量更快稀释了AI占比；结合2026Q1该占比回升至18.2\%，趋势判断需观察后续季度而非单点。

下面进入观察名单与客户侧证据簇。生益电子问询回复首次给出匿名客户逐行收入（R客户2025前9月29.37亿元），证明AI服务器PCB订单正在向具名可追溯的客户编号集中——这是材料端公司第一次出现可进入模型的客户级颗粒度，但仍不能映射到NVIDIA/Google等具名平台。随后的图表5-7至图表5-11构成第二层证据：均为官方披露但链条偏间接（年报具名PCB客户、产品-客户桥、IR刷新），用于校准需求方向而非直接量化EPS。

\begin{exhibitbox}[图表5-6: Shengyi Electronics anonymized AI-server customer bridge]
\footnotesize
\begin{tabularx}{\textwidth}{L{1.8cm}R{1.6cm}R{1.7cm}R{1.8cm}R{1.9cm}>{\raggedright\arraybackslash}X}
\toprule
\textbf{客户} & \textbf{2022 sales} & \textbf{2023 sales} & \textbf{2024 sales} & \textbf{2025 1-9M} & \textbf{解读} \\
\midrule
R客户 & 2.65亿元 & 2.65亿元 & 10.20亿元 & 29.37亿元 & X终端客户指定EMS；AI服务器高附加值PCB量产订单增长。 \\
B客户 & 7.96亿元 & 6.06亿元 & 6.94亿元 & 7.82亿元 & 通信网络设备板、服务器/计算机板。 \\
F客户 & 1.26亿元 & 1.26亿元 & 2.12亿元 & 2.23亿元 & 服务器/计算机板。 \\
C客户 & 1.69亿元 & 2.96亿元 & 2.77亿元 & 1.42亿元（外销口径） & 汽车电子板，非AI主线。 \\
Q客户 & N/A & N/A & N/A & 1.17亿元 & 2025Q3开始明显增长，EMS采购AI服务器高附加值产品。 \\
P客户 & N/A & N/A & N/A & 1.14亿元 & 服务器/计算机板，市场需求带动订单增加。 \\
\bottomrule
\end{tabularx}
\par\vspace{0.6em}
\textit{口径注：}C客户2025 1-9M为外销口径（export sales），其他客户该列均为总额口径（total sales），不可直接横向比较绝对值。
\begin{tabularx}{\textwidth}{L{2.4cm}R{1.7cm}R{1.7cm}R{1.7cm}R{1.7cm}>{\raggedright\arraybackslash}X}
\toprule
\textbf{指标} & \textbf{2022} & \textbf{2023} & \textbf{2024} & \textbf{2025 1-9M} & \textbf{用途} \\
\midrule
外销收入占比 & 41.13\% & 45.51\% & 51.03\% & 63.34\% & 外销和AI服务器客户链权重上升。 \\
服务器/计算机板收入占比 & N/A & N/A & 48.96\% & 68.01\% & 服务器/计算机板成为第一大应用领域。 \\
前五大客户收入占比 & 47.98\% & 47.81\% & 51.28\% & 63.35\% & 客户集中度随AI订单上升。 \\
\bottomrule
\end{tabularx}
\sourcenote{Shengyi Electronics official refinancing inquiry response; customer names exempted, not named NVIDIA/Google revenue}
\end{exhibitbox}

\begin{exhibitbox}[图表5-7: Official named customer-chain evidence from watchlist filings]
\footnotesize
\begin{description}
  \item[南亚新材] 官方年报具名PCB客户包括健鼎科技、奥士康、景旺电子、深南电路、瀚宇博德、生益电子、方正科技、沪电股份、胜宏科技、广东骏亚；终端重点客户认证包括华为、中兴通讯、浪潮、曙光、新华三、HPE、微软。
  \item[高速材料代际] M6--M8产品已批量应用于国内头部算力客户，M9处于NPI导入，M10于2025Q4推出并处于海外核心算力终端认证；NOUYA8U完成华为核心客户准入并规模化量产。
  \item[边界] 年报同时说明因商业秘密，对前五名供应商及客户名称豁免披露。因此该证据能确认客户链和认证，不等于客户/平台收入拆分。
\end{description}
\sourcenote{official named customer-chain evidence; official 2025 annual report of 南亚新材（688519）}
\end{exhibitbox}

5-7给出的是具名PCB客户链与终端认证，但均停留在``产品销售给哪些PCB厂''层面，不携带任何终端平台收入数字。下表5-8转向集中度量化，把客户/供应商占比与折价梯度挂钩。

\begin{exhibitbox}[图表5-8: Watchlist customer and supplier concentration]
\footnotesize
\begin{tabularx}{\textwidth}{L{1.8cm}R{2.0cm}R{1.6cm}R{2.0cm}R{1.6cm}X}
\toprule
\textbf{公司（代码）} & \textbf{Top-5 sales} & \textbf{Sales ratio} & \textbf{Top-5 purchase} & \textbf{Purchase ratio} & \textbf{解读} \\
\midrule
南亚新材（688519） & 18.27亿元 & 34.96\% & 25.99亿元 & 55.34\% & 供应商集中度高，且问询函披露奥士康、方正、深南、景旺等客户。 \\
兴森科技（002436） & 19.64亿元 & 27.29\% & 11.53亿元 & 20.27\% & 客户集中度中等，匿名披露。 \\
大族数控（301200） & 21.74亿元 & 37.65\% & 12.51亿元 & 28.23\% & 第一大客户占27.24\%，设备订单节奏敏感。 \\
芯碁微装（688630） & 5.86亿元 & 41.59\% & 4.22亿元 & 42.24\% & 具名客户含鹏鼎、胜宏、景旺、深南，客户链证据强。 \\
劲拓股份（300400） & 1.47亿元 & 18.72\% & 0.84亿元 & 17.53\% & 集中度较低，订单分散。 \\
鼎泰高科（301377） & 6.44亿元 & 30.04\% & 4.94亿元 & 36.79\% & 耗材链客户/供应商均有集中度风险。 \\
鹏鼎控股（002938） & 350.17亿元 & 89.45\% & 67.29亿元 & 21.28\% & 2025年报口径客户高度集中，第一大客户A占79.65\%，但匿名披露。 \\
\bottomrule
\end{tabularx}
\sourcenote{watchlist customer / supplier concentration summary; Pengding official filing evidence; official 2025 annual reports and inquiry responses}
\end{exhibitbox}

5-8的核心读数是两极分化：鹏鼎前五大客户占比89.45\%（第一大客户A占79.65\%）与兴森27.29\%形成集中度两极，这直接决定折价梯度——集中度越高，单一大客户砍单的敏感度越大，估值折扣也越大。下表5-9聚焦产品-客户桥，把具名客户收入与产品经济性并列，但需特别注意：芯碁微装的设备收入虽含鹏鼎/胜宏/景旺/深南等AI PCB大厂，但设备收入不等于下游AI PCB收入，进入模型时须按传导系数打折。

\begin{exhibitbox}[图表5-9: Watchlist product-customer bridge from official filings]
\footnotesize
\begin{tabularx}{\textwidth}{L{1.8cm}X X X}
\toprule
\textbf{公司（代码）} & \textbf{Product economics} & \textbf{Named customer / validation evidence} & \textbf{Remaining gap} \\
\midrule
南亚新材（688519） & 2025覆铜板收入40.55亿元、毛利率6.96\%；粘结片收入10.72亿元、毛利率31.05\%；高频高速及IC载板材料收入9.25亿元、同比+80.60\%。 & 问询回复列示覆铜板客户奥士康5.86亿元、方正科技2.19亿元、深南电路2.12亿元、江西旭昇1.85亿元、景旺电子1.64亿元；粘结片客户奥士康1.54亿元、方正科技1.39亿元、深南电路1.05亿元、特创集团0.57亿元、景旺电子0.47亿元。M6/M7已稳定量产交付，M8及以下通过部分海外终端认证。 & PCB客户收入强于终端平台收入；无法拆出NVIDIA/Google/海外AI平台贡献。 \\
芯碁微装（688630） & 2025收入14.08亿元、同比+47.61\%；PCB系列收入10.80亿元、毛利率35.59\%；泛半导体系列收入2.33亿元、毛利率54.57\%、同比+112.50\%。 & 年报列示前五客户：鹏鼎1.93亿元/13.68\%、胜宏1.89亿元/13.40\%、景旺1.00亿元/7.07\%、V-Technology 0.57亿元/4.01\%、深南0.48亿元/3.43\%。2026规划聚焦高端PCB、HDI/类载板、WLP/PLP、IC载板曝光设备、4um线宽和CoWoS-L爬坡支持。 & 设备收入不是下游AI平台收入；年报亦披露无重大销售合同。 \\
\bottomrule
\end{tabularx}
\sourcenote{watchlist product / customer bridge; official 2025 annual report inquiry response for 南亚新材（688519）; official 2025 annual report for 芯碁微装（688630）}
\end{exhibitbox}

\begin{exhibitbox}[图表5-10: Additional watchlist segment-customer evidence]
\footnotesize
\begin{tabularx}{\textwidth}{L{1.8cm}X X X}
\toprule
\textbf{公司（代码）} & \textbf{Official product economics} & \textbf{AI-chain / customer evidence} & \textbf{边界} \\
\midrule
兴森科技（002436） & 2025收入71.95亿元；PCB收入48.97亿元、毛利率25.26\%；半导体业务收入19.10亿元，其中IC封装基板16.70亿元、毛利率-16.06\%，半导体测试板2.39亿元、毛利率38.03\%。 & FCBGA项目已完成技术、产能和良率量产准备，样品订单同比大幅增长；但人工、折旧、能源和材料等投入6.62亿元，显著拖累净利。前五客户19.64亿元/27.29\%。 & 客户匿名；IC载板增长不等于AI平台放量，EPS要扣除FCBGA爬坡成本。 \\
大族数控（301200） & 2025 PCB专用设备收入57.73亿元、毛利率35.12\%；钻孔类设备收入41.67亿元、毛利率33.45\%、同比+98.38\%。 & CCD六轴独立机械钻孔机完成下一代AI服务器PCB加工认证，并在多家高多层板龙头企业量产；前五客户21.74亿元/37.65\%，第一大客户15.72亿元/27.24\%。 & 客户匿名，且公司披露无重大销售合同；设备收入不是下游AI PCB收入。 \\
劲拓股份（300400） & 2025收入7.85亿元；电子装联设备收入7.27亿元、毛利率31.68\%、占收入92.59\%；内销占93.02\%，直销占86.41\%。 & 终端应用覆盖消费电子、通信、汽车电子、算力服务器等；AI服务器增长被列为PCBA产业链结构性机会。前五客户1.47亿元/18.72\%。 & 主题相关但链条较远，属于PCBA装联设备代理，不是PCB制造/材料直接收入。 \\
鼎泰高科（301377） & 2025收入21.44亿元；精密刀具收入17.40亿元、毛利率41.67\%；研磨抛光材料1.92亿元、毛利率58.86\%。0.20mm及以下微钻销量占比29.65\%，涂层钻针销量占比39.40\%。 & 高端微钻、高长径比钻针和涂层钻针匹配AI服务器、高端光模块和封装载板加工需求；前五客户6.44亿元/30.04\%。 & 耗材链能验证需求方向，但不能拆分终端AI平台客户或订单金额。 \\
\bottomrule
\end{tabularx}
\sourcenote{watchlist product / customer bridge; official 2025 annual reports for 兴森科技（002436）, 大族数控（301200）, 劲拓股份（300400） and 鼎泰高科（301377）}
\end{exhibitbox}

\begin{exhibitbox}[图表5-11: Watchlist IR technology and customer-validation refresh]
\footnotesize
\begin{description}
  \item[南亚新材] IR记录显示NY6300S进入多客户PCIe Gen5服务器量产，PCIe Gen6完成头部客户测评；AI服务器突破多客户国产GPU方案，材料覆盖M6--M8U并适配UBB/OAM/AI系统；M8材料已获多家国内重要终端认证并小批量生产，M9材料在多家PCB客户测试并推进国内外终端认证。
  \item[芯碁微装] IR记录显示WLP2000晶圆级封装设备精度达2um，已在多家头部客户做量产测试且验证顺利；先进封装中直写光刻在无掩膜、成本、灵活性和大面积高算力芯片曝光效率上具备优势。
  \item[大族数控] IR记录显示CCD六轴独立机械钻孔机搭载3D背钻及钻测一体技术，已通过下一代AI服务器PCB认证并在多家龙头企业量产；AI服务器相关18层以上高多层板、HDI、封装基板需求提升设备门槛。
  \item[边界] 这些IR增强技术/认证证据，但仍未披露客户名称、订单金额或平台收入拆分。
\end{description}
\sourcenote{watchlist official / issuer IR refresh records}
\end{exhibitbox}

最后两图表（图表5-12鹏鼎HKEX、图表5-13客户侧供应链清单）转向客户侧官方证据。其作用需明确界定：客户侧官方清单只能提升供应商关系可信度，不能产生收入数字——价值在于降低名单造假风险，而非量化EPS或平台拆分。读者应将其视为对``供应商关系真实''这一假设的佐证，而非新增的盈利证据。

\begin{exhibitbox}[图表5-12: Pengding HKEX filing customer and technology evidence]
\footnotesize
\paragraph{定位说明} 本表属观察名单/主题储备：证据增强但不改变ch01``暂不进核心''定位，且不替代缺失的UBS原始研报。
\begin{description}
  \item[Market position] HKEX draft filing excerpt cites Frost \& Sullivan: Pengding ranked No. 1 globally by sales revenue in AI and high-performance computing PCB in 1H2025, and No. 1 in high-build-up HDI PCB and 14+ layer high-layer-count MLPCB categories.
  \item[Customer concentration] 2025年前三季度前五大客户收入71.77亿元、占收入50.9\%；第一大客户收入41.72亿元、占29.6\%。鹏鼎Customer E为美国NASDAQ上市的加速计算和AI基础设施公司，2025年前三季度销售5.92亿元、占4.2\%（注：``Customer E''为鹏鼎HKEX文件的匿名客户代码，与胜宏H股招股书中的``Customer E''分属不同公司、不同实体，两份文件的代码体系无对应关系，不可合并解读为同一终端客户）。
  \item[A-share cross-check] 2025年报口径显示前五大客户销售额350.17亿元、占年度销售总额89.45\%；第一大客户A销售额311.82亿元、占79.65\%；客户B为鸿海集团及其控股子公司，销售额24.29亿元、占6.20\%。该口径和HKEX片段口径不同，应并列使用。
  \item[Technology and capacity] 产品覆盖AI加速卡、AI服务器、数据中心交换机、UBB和光模块；具备70层以上MLPCB量产和100层以上MLPCB技术能力，24层6+12+6 HDI和28层8+12+8 HDI量产能力；2025年前三季度HDI利用率91.1\%。
  \item[Official delivery] 2025年汽车/服务器用板及其他用板收入21.19亿元、同比+106.67\%；年报称AI服务器类产品收入较2024年增长超1倍，2026年预计继续保持强劲增长；1.6T光模块业务快速增长，3.2T产品进入研发设计。
  \item[IR validation] 2026年IR记录显示算力客户认证进度顺利，预计2026年是算力直接客户订单导入元年；高阶HDI和HLC产品已打入AI服务器并逐步量产，SLP切入800G/1.6T光模块并将量产出货，3.2T产品进入研发。公司对“英伟达高端板订单”问题明确表示不方便评价单一客户。
  \item[P5W interaction] 深交所互动易/全景问答显示，公司1.6T光模块相关产品已批量供货，3.2T产品与客户开发中；对于Rubin/英伟达、苹果折叠屏等具体客户/产品问题，公司均以商业保密为由不披露。
  \item[母公司映射] 臻鼎2025年报披露营收1,825.22亿新台币、归母净利67.91亿新台币、EPS 6.91元；AI云端、网络与终端应用相关产品收入已接近2025收入70\%。年报还披露Customer X 2025销售1,360.68亿新台币、占74.55\%，但客户匿名，不能映射为Apple/NVIDIA等具名客户。臻鼎官方简报显示，2026Q1 服务器/光模块及其他收入占比10.1\%、IC载板占比9.6\%，合计19.7\%；服务器/光模块收入同比翻倍以上，全年预计翻倍，IC载板全年目标+70\%以上。2025Q3简报进一步称AI服务器收入预计2026年逐步放量、2027年翻倍，淮安iHDI/HLC产能预计2026年底翻倍。
  \item[Apple-side official signal] Apple 2023年官方Newsroom清洁能源公告写明 Avary Holding 于2020年加入 Apple Supplier Clean Energy Program，并启动自身供应商可再生能源倡议。这是客户侧官方供应链项目参与信号，不披露PCB产品、订单、收入、ASP、出货量或平台分配。
  \item[Boundary] 客户代码仍为匿名披露，不能映射为具名NVIDIA/Google/Meta/Microsoft/Apple收入；该证据增强鹏鼎观察名单地位，但不替代缺失的UBS原始研报。
\end{description}
\sourcenote{Pengding HKEX draft filing excerpt, official A-share filings, IR records, P5W/SZSE interaction, Apple Newsroom and parent-company (Zhen Ding) official read-through; not named platform revenue split}
\end{exhibitbox}

5-12给出鹏鼎单公司的客户、技术、产能证据，5-13则把视角扩大到所有核心PCB/CCL供应商在Microsoft/Amazon/Apple/Dell/NVIDIA客户侧的官方供应链清单命中，用于交叉验证供应商关系的真实性。

\begin{exhibitbox}[图表5-13: Customer-side official supplier-list evidence]
\footnotesize
\begin{tabularx}{\textwidth}{L{2.7cm}X X}
\toprule
\textbf{客户侧来源} & \textbf{命中供应商} & \textbf{证据边界} \\
\midrule
Microsoft Top 100 Production Suppliers FY24 & AVARY HOLDING (SHENZHEN), VICTORY GIANT TECHNOLOGY (HUIZHOU), HANNSTAR BOARD, TRIPOD TECHNOLOGY, UNIMICRON TECHNOLOGY (KUNSHAN), SUZHOU DONGSHAN PRECISION, SAMSUNG ELECTRO MECHANICS, MEKTEC。 & Microsoft官方商用硬件Top 100供应商清单，能确认客户侧具名供应商关系；不披露具体产品、PCB收入、AI/云平台归属、ASP、出货或订单金额。 \\
Amazon 2024 Sustainability Report & 报告称供应商清单含近2,300个成品与组件供应商，并共享至Open Supply Hub。 & 当前报告正文未列出PCB供应商行；OS Hub匿名API需要认证。 \\
Open Supply Hub facility API / front-end & Victory Giant、Avary 初始设施页显示 Amazon 贡献者；扩展检索新增 Tripod 3个相关设施（Amazon/Apple/Dell/Samsung/AWS，含 Parts/Components、Other direct material suppliers、Semiconductor Manufacturing、Finished goods 等粗粒度字段）、Unimicron 6个设施（Apple/AWS）、Dongshan 3个设施（Apple）、Meiko 5个设施（Amazon/Samsung/gBizINFO）、Mektec 1个设施（Amazon/Apple）、Shennan Circuits USA 1个SBA行。 & 客户侧网络贡献者和设施元数据证据；不披露具体产品、PCB收入、AWS/AI平台归属、ASP、出货、订单金额、毛利或EPS底层假设。 \\
Upstream supplier-list files & Apple 2018/FY2020名单命中 Dongshan、Tripod、Unimicron、Samsung Electro-Mechanics、Zhen Ding；Dell FY2025名单命中 Tripod、HannStar/Gold Circuit类零组件行，Tripod为 Parts / Components、Other direct material suppliers；Samsung供应商名单命中 Meiko、Samsung Electro-Mechanics、Tripod、Korea Circuit、Ibiden、Daeduck；AWS认证站点命中Avary/宏启胜、Qing Ding/臻鼎系、Tripod和Victory Giant。 & 上游客户/认证清单能回溯部分OSH贡献者来源并提高关系可信度；仍不披露具体客户平台产品、收入、ASP、出货、订单金额、毛利或EPS假设。 \\
Customer annual reports / Form SD & Apple披露制造采购义务562亿美元；Alphabet披露1491亿美元承诺，主要为技术基础设施和库存订单；Microsoft披露1099.53亿美元采购承诺，主要与数据中心相关；Dell披露188亿美元采购义务、AI-optimized servers收入246.83亿美元；NVIDIA Form SD披露fabless/contract manufacturing、memory/substrates/components、164个direct suppliers和246个3TG processing facilities。 & 这些一手文件强化AI基础设施和硬件供应链需求、采购承诺和供应商尽调边界，但仍不披露Victory Giant/Avary/Tripod等PCB/CCL/载板供应商的客户平台收入、ASP、出货、订单金额、毛利或EPS变量。 \\
NVIDIA Q2 FY2026 IR release & 未命中Victory Giant / PCB / substrate / ODM qualified相关具名表述。 & NVIDIA官方发布未确认社媒中的具名供应商说法。 \\
\bottomrule
\end{tabularx}
\sourcenote{customer-side official supply-chain evidence: Microsoft Top 100 supplier list, Open Supply Hub facility data, upstream supplier-list files, customer annual reports / Form SD, and customer purchase-commitment matrix; not revenue split}
\end{exhibitbox}

\begin{exhibitbox}[图表5-14: Official filing segment evidence]
\footnotesize
\begin{tabularx}{\textwidth}{L{1.6cm}X X}
\toprule
\textbf{公司} & \textbf{硬证据} & \textbf{缺口} \\
\midrule
沪电股份 & 官方年报披露2025年PCB收入约181.43亿元，毛利率约36.91\%；数通PCB收入约146.56亿元，高速交换机/路由81.69亿元，AI服务器/HPC 30.06亿元。 & 未披露客户名称、NVIDIA/ASIC拆分、单机价值量和平台订单。 \\
深南电路 & 官方年报确认PCB、电子装联、封装基板的3-In-One布局；券商全文给出PCB收入143.59亿元、载板收入41.48亿元。 & 未披露AI服务器订单金额和FC-BGA客户结构。 \\
生益科技 & 官方年报、2025H1、2025Q3和生益电子再融资问询回复已归档；官方文件给出分部收入、AI HDI/高多层算力板扩产、5--6阶HDI客户认证和订单储备。 & M8/M9收入占比、具名认证客户、涨价幅度和客户收入仍缺。 \\
华正新材 & 官方修订版年报和2026高等级覆铜板可行性报告已归档；P5W/上交所互动补充CBF小批量订单与类BT/BT应用边界；浙商深度和年报点评补充CBF/BT材料、2025E--2028E预测线和三表附录。 & CBF/BT收入、客户/平台收入拆分和项目级底层假设仍缺。 \\
\bottomrule
\end{tabularx}
\sourcenote{official filing segment summary; official 2025 annual reports retrieved from CNInfo/SZSE/SSE}
\end{exhibitbox}

综合上述观察名单与客户侧证据，可按证据密度归为三簇，便于判断每条证据在盈利模型中的权重。\textbf{第一簇：核心公司官方分部证据密度最高。}沪电股份的数通PCB/高速交换机/AI服务器分部、深南电路的PCB/载板分部、生益科技的覆铜板/PCB分部均为公司年报披露，可进入收入与毛利率模型；但其共同边界是\emph{不披露具名平台（NVIDIA/Google ASIC/光模块）收入}，因此任何“AI相关占比”若非官方分部口径，只能作为待验证情景变量，不可直接写入EPS。\textbf{第二簇：观察名单证据以官方披露为主、但链条偏间接。}南亚新材/芯碁微装/兴森/大族数控/鼎泰高科的年报与再融资问询回复能确认客户链与认证状态（如M6--M8量产、CCD钻孔机通过AI服务器PCB认证、芯碁微装前五客户含鹏鼎/胜宏/景旺/深南），但这些是材料/设备/耗材收入，\emph{不等于下游AI平台PCB收入}，进入模型时须按下游传导系数打折并标注证据层级。\textbf{第三簇：客户侧官方供应链证据提升关系可信度，但不构成收入。}Microsoft Top 100供应商、Open Supply Hub设施元数据、Apple/Alphabet/Microsoft/Dell/NVIDIA的Form SD与采购承诺，能回溯Avary/Victory Giant/Tripod/Unimicron等供应商的客户侧关系，但同样不披露PCB收入、ASP、出货或订单金额——其价值在于\emph{佐证供应商名单真实性}，而非量化贡献。三簇共同决定了本章的核心边界：平台链条是可验证的需求来源，但具名客户/平台收入拆分一律未公开披露。

\section{扩充标的客户链分层}

第二批14只扩充标的的客户链证据密度整体低于核心6家，需按层级区别对待。Core Candidate 3家（东山精密、景旺电子、芯碁微装）已有可验证的客户链线索，但证据强度弱于核心公司；Demand Indicator（工业富联、中际旭创）提供反向视角——它们是PCB/CCL的需求方，客户链需从"谁采购它们的产品"角度理解；其余Watchlist标的客户链证据参差不齐，部分仅有行业层面间接证据，部分已公开澄清无AI客户。整体原则：不编造客户名称，所有客户链陈述均标注证据强度（强/中/弱），缺口明确标记。

\subsection{Core Candidate 客户链}

Core Candidate 三家定位于核心公司与观察名单之间：客户链方向可确认、AI 业务敞口有证据，但颗粒度尚未达到核心公司的分部披露水平。

\begin{exhibitbox}[图表5-15: Core Candidate 客户链证据]
\footnotesize
\begin{description}
  \item[东山精密（002384）—— 证据强度：\textbf{中}]
    PCB 百强第 2，AI 服务器 PCB + 光模块/光芯片垂直整合（索尔思光电 2025 年 10 月并表）。客户侧：英特尔 AI-server PCB 核心供应商为市场公认口径（公司文件未直接具名）；Google ASIC 关联仅见于第三方评论，未在公司年报客户列表中验证。客户类别覆盖全球 AI 技术解决方案供应商、大型云服务商、数据中心设备 OEM、服务器厂商等（H 股招股书级别描述，客户名称匿名）。2026Q1 营收 131.4 亿元（YoY +52.72\%）、归母净利 11.10 亿元（+143.47\%），业绩高速增长验证需求侧拉动，但客户级收入拆分不可得。缺口：具名客户（NVIDIA/Google/英特尔）收入占比、AI PCB 与光模块收入拆分、单机价值量。
  \item[景旺电子（603228）—— 证据强度：\textbf{中}]
    全球主要 AI 服务器 PCB 制造商，汽车电子 PCB 全球排名第一。客户侧：软板与软硬结合板已打入 AI 服务器/数据中心，在英伟达供应链具稀缺卡位（厂商自述 + 卖方转引，未获 NVIDIA 客户方公开背书）；Rubin 架构相关背板产品推进中。1.6T 光模块 PCB 已实现产业化作业并开始出货（公司 2026-03-30 投资者关系活动记录表，高置信）；800G 已批量稳定供货头部客户（客户匿名）。泰国巴真府基地 2026 年投产，海外客户认证推进中。2025 年营收 153.08 亿元、归母净利 12.31 亿元；2026Q1 归母 -28.37\%，利润端承压。缺口：NVIDIA/Rubin 具名客户确认、AI 服务器 PCB 收入占比、1.6T 光模块 PCB 出货量级与单价。
  \item[芯碁微装（688630）—— 证据强度：\textbf{强}]
    全球 PCB 直写光刻设备市占率第三（约 18.8\%，卖方 + 招股书口径），客户链证据为扩充标的中最强。2025 年报前五具名客户：鹏鼎 1.93 亿元（13.68\%）、胜宏 1.89 亿元（13.40\%）、景旺 1.00 亿元（7.07\%）、V-Technology 0.57 亿元（4.01\%）、深南 0.48 亿元（3.43\%）。五家客户均为 AI PCB 核心供应商，设备采购与 AI 扩产周期直接联动。2026Q1 营收 5.15 亿元（+112.48\%）、归母净利 1.08 亿元（+108.98\%）、经营性现金流 +1.65 亿元（+625.51\%），三表同步跳升验证扩产传导。边界：设备收入不等于下游 AI 平台收入，进入盈利模型须按下游传导系数打折。缺口：CO\textsubscript{2} 激光钻孔客户逐项具名、先进封装 WLP/PLP 客户收入拆分。
\end{description}
\sourcenote{Core Candidate customer-chain evidence; official 2025 annual reports; issuer IR records; broker report and industry league table citations; customer names verified against annual report top-5 disclosures where available}
\end{exhibitbox}

Core Candidate 三家的共同特征是客户链方向已可确认、但颗粒度不足。东山精密和景旺电子的 AI 服务器 PCB 客户均停留在"厂商自述 + 市场公认"层面，缺具名客户收入拆分；芯碁微装的设备客户虽全部具名，但属"铲子股"传导链，AI 弹性需经下游 PCB 厂的产能利用率和订单转化率二次校准。

\subsection{Demand Indicator 客户链（反向视角）}

工业富联与中际旭创定位为需求侧锚点，客户链视角与 PCB/CCL 标的相反：不是"它们是谁的客户"，而是"谁是它们的客户"。这一反向视角的价值在于验证 AI 算力需求是否真实落地到下游系统级厂商的收入和订单中，从而反向支撑上游 PCB/CCL 的需求逻辑。

\begin{exhibitbox}[图表5-16: Demand Indicator 客户链（反向视角）]
\footnotesize
\begin{description}
  \item[工业富联（601138）—— 证据强度：\textbf{中}]
    下游需求侧锚，AI 服务器 + 高速交换机 + CPO 三大业务均指向 AI 算力基础设施。客户侧：与全球顶尖科技公司（NVIDIA 及北美超大规模云厂商）建立深度绑定，CPO 领域实现 L11 全层级绑定（公开搜索聚合，非一手归档）。800G 及以上交换机 2025H1 营收已达 2024 全年近 3 倍；GB200/GB300 出货顺畅；CPO 全光交换机样机已生产，1.6T 交换机 2026 年放量。2026Q1 营收 2,510.78 亿元（+56.52\%）、归母净利 105.95 亿元（+102.55\%）、经营性现金流 250.24 亿元（+1,826\%），三表同步爆发验证需求强度。边界：客户名称未在公司年报中逐项具名；为需求指标股，非 PCB-CCL 链内投资标的。
  \item[中际旭创（300308）—— 证据强度：\textbf{中}]
    全球 800G/1.6T 光模块龙头，需求侧拉货锚标。客户侧：下游为北美云厂商与 AI 服务器厂商；谷歌 2026 年 1.6T 采购占全球预测需求 35--40\%（产业链调研口径，非客户公告）；GB300 采用 1.6T 模块比例由 1:2.5 提升至 1:5（产业链调研具名）。2026 年 1.6T 产能规划约 1,500 万只，预计出货约 1,200 万只（公司口径 + 卖方）。2026Q1 营收 194.96 亿元（+192.12\%）、归母净利 57.35 亿元（+262.28\%）。边界：客户名称未具名披露；PCB/CCL 供应商名单未披露，生益/华正/胜宏/沪电是否进入其 BOM 属待 IR 验证范畴；为光模块下游需求锚，非 PCB 制造标的。
\end{description}
\sourcenote{demand-side indicator customer evidence (reverse perspective); official filings and issuer IR records; industry research and supply-chain channel checks; not named customer revenue split; demand indicators, not PCB-CCL chain investable names}
\end{exhibitbox}

两家需求锚的核心读数是：AI 算力需求已从 CAPEX 指引转化为系统级厂商的真实收入和现金流。工业富联的 800G 交换机 3 倍增长、中际旭创的 1.6T 出货量级，共同构成了从"北美云厂 AI capex"到"中国高端 PCB/CCL 订单"之间的需求传导中继站——但这一传导仍为链条级逻辑，不等同于任何单家 PCB/CCL 公司的具名客户收入。

\subsection{Watchlist 扩充标的客户链}

Watchlist 扩充标的共 8 家（铜冠铜箔、东材科技、宏和科技、金安国纪、中材科技、中国巨石、宏昌电子、容大感光），客户链证据差异较大。下表按证据强度排序，明确标注每只标的的客户结构、AI 链关联与缺口。

\begin{exhibitbox}[图表5-17: Watchlist 扩充标的客户链证据（按强度排序）]
\footnotesize
\begin{tabularx}{\textwidth}{L{2.2cm}>{\raggedright\arraybackslash}p{1.0cm} >{\raggedright\arraybackslash}X >{\raggedright\arraybackslash}X}
\toprule
\textbf{公司（代码）} & \textbf{强度} & \textbf{客户链证据} & \textbf{缺口 / 边界} \\
\midrule
铜冠铜箔（301217） & \textbf{中} & HVLP4 铜箔已批量供应 AI 服务器及光模块领域；客户含生益科技、胜宏科技等头部 PCB 板厂（厂商 + 卖方一致）。PCB 铜箔 2025 年收入 37.04 亿元（占比 55\%），是 A 股唯一 PCB 铜箔收入过半的纯标的。HVLP5 代铜箔关键性能指标研发突破，中试收尾。 & 客户收入占比、HVLP4 单价与毛利率贡献未披露；生益/胜宏方未公开背书；HVLP5 与 IC 载体铜箔尚未量产。 \\
东材科技（601208） & \textbf{中} & 高速树脂国内绝对龙头（BMI/苯并噁嗪/碳氢国内市占第一），通过台光、台耀、生益科技、斗山等全球主流 CCL 厂商进入供应链，最终覆盖 NVIDIA H100/H200/GB300、华为昇腾 910B、Apple 等终端。M10 正交背板方案已通过 NVIDIA 终端认可，进入下一阶段验证。2026Q1 营收 +27\%、归母 +103\%。 & 终端 NVIDIA 认可为 CCL 厂间接传导，非直接客户订单；具名 CCL 客户为公司口径，未获客户方交叉验证；AI 树脂收入占比未拆分。 \\
宏昌电子（603002） & \textbf{弱-中} & 国内电子级环氧树脂产能第一，覆盖 M2--M9 全系列。环氧树脂直接客户包括瀚宇博德、健鼎科技等大型 PCB/CCL 厂商。子公司无锡宏仁高速覆铜板（GA-686 / GA-686N）已通过 Intel、AMD 评估认证并取得小量订单。 & Intel/AMD 认证为公司口径，未获客户方确认；高速 CCL 仍处小量订单阶段；2026Q1 归母净利仅 46.85 万元（-92.74\%），"放量但亏利润"剪刀差显著。 \\
宏和科技（603256） & \textbf{弱-中} & 国内少数完成低介电高端电子布客户认证并稳定出货的厂商，产品最终应用于 AI 服务器高速 CCL。第三代 Low-Dk 与石英纤维布为升级方向。2025 全年电子布/电子纱营收约 11.71 亿元、归母净利 +786\%。 & "完成认证 + 稳定出货"为公司自述，未取得生益/南亚/台光等具体 CCL 客户采购合同或认证函原文；AI 电子布收入占比未在年报中分项披露。 \\
中材科技（002080） & \textbf{弱} & 泰山玻纤已覆盖低介电一代、二代、低膨胀布、超低损耗低介电布（Q 布）全品类，公司口径称"均已完成国内外头部客户认证"。2025 年特种纤维布收入 6.28 亿元、占总营收 2.08\%。定增近 45 亿元投向 3,500 万米低介电 + 2,400 万米超低损耗 Q 布项目。 & "头部客户认证"为公司自述，第三方交叉验证不足；AI 电子布收入占比仅 2.08\%，当期贡献极小；2026Q1 经营性现金流 -20.14 亿元（-599\%）。 \\
容大感光（300576） & \textbf{弱} & PCB 光刻胶国产替代标的，湿膜光刻胶国内市占率约 50\%（媒体口径）。部分 AI 算力 PCB 光刻胶型号已量产，正推进头部客户验证（公司 IR 口径）。传导链：AI 服务器 PCB → 高端感光油墨/干膜 → 容大感光。 & AI 光刻胶收入未在年报分部中具名披露；客户名称未公开；2026Q1 归母 -22.43\%、扣非 -23.73\%、经营性现金流仍为负，三表不支持 AI 放量结论。 \\
金安国纪（002636） & \textbf{弱} & CCL 行业前三、第二梯队龙头，主打中低端 FR-4 覆铜板，客户为 PCB 厂商（匿名披露）。间接受益于 CCL 全行业涨价周期，2026Q1 归母 +764\%。 & 公司 2026-06-15/16 两次澄清：未与英伟达、华为合作，未做高速 M6/M9，未发现 AI 服务器应用。AI 敞口 ≈ 0，估值中 AI 概念溢价部分无基本面支撑。 \\
中国巨石（600176） & \textbf{弱} & 全球玻纤龙头（市占率约 24\%）、电子纱/电子布全球份额约 23\%。普通电子布为 CCL 基材，受益于电子布行业涨价周期。2026-05 公告拟投建 5 万吨电子纱 + 3.2 亿米电子布项目（44.31 亿元）。 & 公司 2026-06-17 主动澄清：低介电等特种电子布产品未形成订单及收入。AI 链敞口当前严格为 0；利润弹性来自普通电子布/粗纱周期，非 AI 链 α。 \\
\bottomrule
\end{tabularx}
\sourcenote{watchlist expanded-name customer-chain evidence; official annual reports, issuer clarifications, IR records, and broker report citations; evidence strength self-rated based on named customer disclosure vs. anonymous / company-self-reported claims}
\end{exhibitbox}

八家 Watchlist 扩充标的的客户链呈现明显梯度：铜冠铜箔与东材科技因直接面向生益/胜宏/台光等 AI CCL 核心客户而证据相对较强，但均为"上游材料 → CCL 厂 → AI 终端"的间接传导，AI 收入占比不可得；宏昌电子、宏和科技有一定客户线索但缺乏第三方交叉验证；中材科技、容大感光仅有公司口径陈述；金安国纪、中国巨石则已公开澄清无 AI 客户/订单，其 AI 关联仅为市场叙事。

\subsection{建滔积层板（港股）}

建滔积层板（01888.HK）为全球 CCL 龙头之一，垂直一体化程度仅次于台光电，是 A 股 CCL 标的的全球定价锚与量价风向标。客户链证据较少：客户为全球 PCB 厂商（分散、未逐项具名），high-M-grade（M6/M8/M9 级）超低损耗 CCL 为 AI 服务器基材主力供应方之一（行业第一梯队，与台光、生益、斗山、松下同列）。2025H2--2026 累计提价 >50\%，是 CCL 涨价周期的最强外生信号。\textbf{证据强度：弱-中}——行业地位强但具名客户披露少，且港股信息披露标准与 A 股不同。客户结构以年报分部为准，不做具名推断。

\section{与顶级研究的差距}

\begin{sourcequalitybox}[Main content gap]
当前报告能证明``平台链条是主要需求来源''，但尚不能证明``每条平台链条对每家公司EPS贡献多少''。这是与顶级全球硬件研报相比最大的内容差距，也是下一轮调研的最高优先级。
\end{sourcequalitybox}

因此，本报告的投资建议不应被理解为“AI PCB全链条买入”。更准确的理解是：对证据密度高、收入和毛利已经开始兑现的公司，维持核心跟踪；对客户和产品弹性更大的公司，等待Q2/中报验证；对设备、耗材和材料期权，要求更高的安全边际和更明确的订单证据。

\section{投资者交流平台证据}

\begin{exhibitbox}[图表5-18: IR records improve platform-chain evidence]
\footnotesize
\begin{description}
  \item[沪电股份（002463）] 2026年6月IR称AI训练/推理和CSP资本开支推动AI服务器、高速网络交换机部署；泰国数据通讯事业部已有超70\%海外客户认证，2026Q1产能利用率超90\%，加速AI服务器与高速网络中高端产品批量导入。缺口：未披露终端客户名称或NVIDIA/ASIC收入。
  \item[胜宏科技（300476）] 2026年6月IR引用Prismark第三方市场研究数据：AI服务器PCB 2024--2029 CAGR 18.7\%，AI服务器HDI CAGR 29.6\%（市场规模增速为第三方测算，非公司指引）；公司具备100层以上高多层PCB、10阶30层HDI、16层任意互联HDI；ASIC客户进展顺利，GPU加速卡和TPU配套板供应规模扩大；1.6T光模块mSAP产线在手订单饱满。缺口：未披露具体客户或ASIC/NVIDIA平台收入。
  \item[深南电路（002916）] 2026年6月IR称PCB业务受益于AI算力基础设施硬件需求增长，综合产能利用率处于高位；封装基板因处理器芯片和存储类需求拉动；FC-BGA 22层及以下产品量产、24层以上研发推进；2026资本开支聚焦PCB和封装基板。缺口：未披露平台/客户订单金额。
  \item[生益科技（600183）] 上证e互动官方IR记录显示，AI服务器对低损耗CCL材料提出更高要求；公司正同国内外终端围绕GPU和AI项目开发合作，并已有产品批量供应；泰国CCL/prepreg基地投资约14亿元。缺口：未披露M8/M9收入占比或具名客户认证。
\end{description}
\sourcenote{IR platform-chain summary; official IR refresh (2026年6月); official IR PDFs (2026年6月) retrieved from issuer IR portals}
\end{exhibitbox}

\section{客户集中度}

客户集中度决定盈利弹性和估值折扣的两个方向：AI订单集中推高短期弹性，但任一大客户节奏变化会同步压低利润和倍数。下表用官方年报口径披露前五大客户占比，并用匿名逐行明细给出最大客户权重。集中度排序直接映射估值处理——沪电（53\%）和胜宏（42\%）的高集中度解释其高盈利弹性，但也意味着任一大客户砍单会同步打击收入和倍数，应给集中度折扣而非纯增长溢价；深南（18\%）的分散组合意味着更稳但更低的弹性，更适合作为防御性配置而非AI高弹性头寸。

\begin{exhibitbox}[图表5-19: Official top-5 customer concentration]
\small
\begin{tabularx}{\textwidth}{L{1.8cm}R{2.4cm}R{2.0cm}X}
\toprule
\textbf{公司（代码）} & \textbf{Top-5 sales} & \textbf{比例} & \textbf{解读} \\
\midrule
沪电股份（002463） & 101.03亿元 & 53.32\% & 客户集中度高；按终端客户口径约53.89\%。 \\
胜宏科技（300476） & 80.98亿元 & 41.98\% & 高集中度，解释其对国际头部客户订单节奏敏感。 \\
深南电路（002916） & 43.56亿元 & 18.42\% & 客户集中度相对较低，组合更分散。 \\
生益科技（600183） & 82.80亿元 & 29.71\% & 材料龙头客户集中度中等偏高。 \\
华正新材（603186） & 11.54亿元 & 26.42\% & 客户集中度中等，客户名称未披露。 \\
\bottomrule
\end{tabularx}
\sourcenote{customer concentration summary; official annual reports}
\end{exhibitbox}

下表进一步给出最大客户的逐行金额与占比，并标注A股年报与H股招股书两套口径的差异，用于判断单一大客户砍单的敏感度。

\begin{exhibitbox}[图表5-20: Anonymous top-customer detail from filings]
\footnotesize
\begin{tabularx}{\textwidth}{L{1.8cm}X X}
\toprule
\textbf{公司（代码）} & \textbf{Top customer rows} & \textbf{解释} \\
\midrule
沪电股份（002463） & 客户一27.02亿元/14.26\%；客户二21.54亿元/11.37\%；客户三20.68亿元/10.92\%；客户四18.04亿元/9.52\%；客户五13.74亿元/7.25\%。 & 高度集中，且发行人按匿名客户披露。 \\
胜宏科技（300476） & 第一名28.87亿元/14.97\%；第二名24.23亿元/12.56\%；第三名18.37亿元/9.52\%；第四名5.20亿元/2.70\%；第五名4.31亿元/2.23\%。 & 前三客户贡献明显，订单节奏敏感。注：本表为A股年报口径，前五大客户合计80.98亿元/41.98\%、最大客户28.87亿元/14.97\%；H股招股书口径2025年前五大客户51.0\%、最大客户29.7\%（Customer A 57.38亿元）。两套口径差异主要源于分部/合并及客户划分标准不同，应并列使用。 \\
深南电路（002916） & 第一名17.65亿元/7.46\%；第二名7.28亿元/3.08\%；第三名6.92亿元/2.93\%；第四名5.95亿元/2.52\%；第五名5.76亿元/2.43\%。 & 客户集中度相对较低。 \\
生益科技（600183） & 前五客户合计82.80亿元/29.71\%。 & 抽取文本未可靠给出逐行客户明细。 \\
华正新材（603186） & 前五客户合计11.54亿元/26.42\%。 & 原文存在分组单位，逐行结构不稳定。 \\
\bottomrule
\end{tabularx}
\sourcenote{top customer / supplier detail from filings; official annual reports}
\end{exhibitbox}

供应商集中度是上游材料/基材供应风险的同源指标，与客户集中度一起决定公司对单一下游周期的暴露。华正新材前五供应商占比高达46\%、生益科技亦达27\%，意味着材料代际切换或上游产能波动会经由成本端放大盈利扰动；这一维度在客户集中度之外补充了第二层风险来源。

\begin{exhibitbox}[图表5-21: Anonymous top-supplier detail from filings]
\footnotesize
\begin{tabularx}{\textwidth}{L{1.8cm}X X}
\toprule
\textbf{公司（代码）} & \textbf{Top supplier rows} & \textbf{解释} \\
\midrule
沪电股份（002463） & 供应商一31.56亿元/22.28\%；供应商二9.72亿元/6.86\%；供应商三7.23亿元/5.11\%；供应商四5.93亿元/4.18\%。 & 上游集中度高，材料供应风险重要。 \\
胜宏科技（300476） & 第一名16.50亿元/9.60\%；第二名12.56亿元/7.31\%；第三名11.65亿元/6.77\%；第四名6.93亿元/4.03\%；第五名6.55亿元/3.81\%。 & 上游集中度中等偏高。 \\
深南电路（002916） & 第一名12.25亿元/8.69\%；第二名9.97亿元/7.08\%；第三名7.70亿元/5.46\%；第四名7.19亿元/5.10\%；第五名5.29亿元/3.75\%。 & 材料供应链稳定性影响PCB/载板交付。 \\
生益科技（600183） & 前五供应商合计50.88亿元/27.45\%。 & 材料企业仍有供应集中风险。 \\
华正新材（603186） & 前五供应商合计16.20亿元/46.12\%。 & 供应商集中度高。 \\
\bottomrule
\end{tabularx}
\sourcenote{top customer / supplier detail from filings; official annual reports}
\end{exhibitbox}

综合三表，客户与供应商集中度共同界定了本报告五家核心公司的风险结构：沪电和胜宏是高弹性-高集中度资产（适合进攻但需折价），深南是分散组合的稳健资产，生益和华正集中度中等、风险更多来自材料代际而非客户。这一结构支撑ch01的核心-观察名单划分——沪电、胜宏、深南因弹性或稳健进入核心，生益和华正因材料代际暴露留在观察——并决定了下文估值不能对五家公司一刀切套用AI硬件倍数，必须按集中度折扣和风险来源分别定价。
